% Part 5: gravitational decoherence, the virial partition and the persistence of classical structure.
% Carried line by line from: Gravitational Decoherence, the Virial Partition, and the Emergence of Classical Structure (Mar 2026; abstract, sections 1,
% 2, 3, 4, 6, Fig. 1). Its observational section 5 is carried in Chapters ch:virial and ch:virial_tests; its halo-survival prediction is tested in
% Chapter ch:satellites. The acknowledgments are not carried (they name private correspondents).
% Corrections applied: errata V1, V9, V10, V11, V16, V21, V22, V30; VIRIAL_CHECK.md items 1, 9-12, 19, 20, 28; ledger note 12(e)-(f).
% Numbers: docs/verification/scripts/verify_virial_papers.py (sections E, G, K).

\chapter{Gravitational decoherence, the virial partition and the persistence of classical structure}\label{ch:virial_decoherence}

\section*{Overview}
The question of why time has a direction, and why classical structure persists rather than dissolving back into quantum superposition, has a thermodynamic
answer that connects to cosmological observables. This chapter proposes that both questions share a common physical mechanism: the irreversible partition of
gravitational binding energy into a geometric channel and an informational channel, governed at every scale by the virial theorem for $1/r$ potentials.
\conjecture

The geometric half of each virialization event is deposited into spacetime curvature. The informational half --- the Landauer cost of the
quantum-to-classical transition --- must be written to an available encoding surface. At galactic scales and above, a local black hole horizon is proposed
as the encoding vault for it; where no such vault can form, the partition is proposed not to close, and the structure to disperse. The accumulated
informational cost of 13.8 billion years of such transitions is measured by the activation function $E(a)=\exp(1-1/a)$, derived from horizon
thermodynamics via the Jacobson--Cai--Kim formalism~\cite{Jacobson1995,CaiKim2005}. This function is monotonically increasing by construction, irreversible,
negligible before structure forms, and asymptoting to $e$ without arriving --- an arrow of time from the irreversibility of each record, not a statistical
artefact of initial conditions. \interp\ The cosmological consequences are tested against Planck 2018 with the coupling fixed at its prediction, consistent
with Planck ($\Delta\chi^2=+0.54$; Chapter~\ref{ch:virial}), and the decisive test is the growth of structure measured by Euclid and DESI
(Chapter~\ref{ch:virial_tests}).

\section{Two faces of the problem of time}\label{sec:vd_intro}
The problem of time has two faces that are worth distinguishing. The first is the direction of time --- why processes are irreversible, why entropy
increases, why the past differs from the future. The second is the persistence of classical structure --- why quantum superpositions give way to definite
outcomes, and why those outcomes cohere into stable macroscopic objects rather than dissolving back into the quantum background.

These two problems are usually treated separately. The direction of time is assigned to statistical mechanics and the second law. The persistence of
classical structure is assigned to decoherence theory~\cite{Zurek2003}. Neither programme has produced a physical mechanism that connects the irreversibility
of individual quantum-to-classical transitions to the large-scale structure of the universe or to measurable cosmological observables.

IAM proposes that both problems have a common physical origin: the virial partition of gravitational binding energy at every scale. For any bound system
governed by a $1/r$ potential, the virial theorem partitions energy exactly in half between a geometric channel and a kinetic channel. In IAM, the kinetic half
is identified as the Landauer cost~\cite{Landauer1961} of the quantum-to-classical transition --- the thermodynamic price paid each time a decoherence event
writes classical information permanently to an encoding surface (Chapter~\ref{ch:virial_identity}).

This identification has a precise consequence. For the arrow of time: each decoherence event is irreversible by the second law, and its Landauer cost is
written permanently on the cosmological apparent horizon~\cite{Jacobson1995,CaiKim2005,GibbonsHawking1977}. The accumulated ledger of all such events
constitutes a thermodynamically necessary arrow of time --- not imposed, but following from the irreversibility of each individual transaction. For the
persistence of structure: the informational half of each virialization event requires an encoding surface. Where a black hole horizon can form, the Landauer
cost is proposed to be absorbed locally and the virial partition to close. Where no local surface can form, the partition is proposed not to close, and the
structure to disperse. This would provide a physical account of why some candidate structures survive as bound systems and others do not. \conjecture

\section{The virial partition as the quantum--classical connector}\label{sec:vd_connector}
For any potential of the form $V\propto r^{-1}$, the virial theorem states $2K=|V|$, $K=\tfrac12|V|$ (Eq.~\eqref{eq:vl_virial}). This is exact for all $1/r$
potentials --- gravitational and Coulomb alike. It is a mathematical theorem, not an approximation, governing every bound system from hydrogen atoms to
galaxy clusters. \derived

Within IAM, the physical interpretation is precise. The potential half $|V|/2$ is deposited into spacetime curvature --- the geometric channel that general
relativity describes. The kinetic half $K=|V|/2$ is the decoherence energy: the Landauer cost of converting quantum superpositions to classical outcomes at
each virialization event. Every time a dark matter halo virializes, half its gravitational binding energy goes into irreversible information production,
written on the nearest available encoding surface at thermodynamic cost $k_BT\ln2$ per bit~\cite{Landauer1961}. \interp\ This identification gives the
cosmological coupling constant immediately: the fraction of matter energy density participating in decoherence per Hubble time is $\beta_m=\Omega_m/2$,
derived from the virial theorem with zero free parameters and held fixed in every Planck chain (Chapter~\ref{ch:virial}).

The $1/2$ partition is not specific to cosmological scales. It is verified across every physical system governed by a $1/r$ potential, from hydrogen atoms
to clusters of galaxies, and it fixes the coupling at the cosmic horizon --- 37 decades in scale, same $1/2$, same mechanism (Chapter~\ref{ch:virial_law}).
This universality is the key point. The fraction of energy going into quantum-to-classical transitions is $1/2$ at every scale not because of a special
property of any particular system, but because every $1/r$ potential obeys the same theorem. The virial theorem is the universal scaling law connecting local
quantum-to-classical transitions to cosmological structure. \interp

\section{Accumulated decoherence and the arrow of time}\label{sec:vd_arrow}
\paragraph{Three temporal quantities.} It is worth distinguishing among three quantities that physical theory often conflates (Chapter~\ref{ch:time}).
\emph{Coordinate time} is a label on the spacetime manifold. It assigns a number to each event without requiring anything to happen. It has no preferred
direction and accumulates nothing; in the relational tradition it is redundant as a fundamental concept (Barbour \& Bertotti 1982~\cite{BarbourBertotti1982}).
\emph{Proper time} $\tau$ is the geometric path length along a worldline, $\tau=\int\sqrt{-g_{\mu\nu}dx^\mu dx^\nu}/c$. It is a property of the spacetime
manifold measured along a particular path. It distinguishes clocks from one another but does not, by itself, distinguish past from future. For photons on
null geodesics, $d\tau=0$ identically. \emph{Accumulated decoherence} is a third quantity. It is thermodynamic, not geometric. It counts the accumulated
irreversible quantum-to-classical transitions on a worldline --- the decoherence events that write classical information permanently on the cosmic horizon at
Landauer cost~\cite{Landauer1961} $k_BT_H\ln2$ per bit, where $T_H=\hbar H/2\pi k_B$ is the Gibbons--Hawking temperature~\cite{GibbonsHawking1977} ($2.65\times10^{-30}$\,K
today at $H_0=67.16$ \calc). It requires massive particles on timelike worldlines. It cannot accumulate on null geodesics. It has a beginning: the electroweak symmetry
breaking, after which stable massive particles first exist (Chapter~\ref{ch:electroweak}). The irreversibility itself does not need CP violation:
decoherence and entropy production are irreversible for any interaction. CP violation enters through the baryon asymmetry (Sakharov
1967~\cite{Sakharov1967}), not through the arrow of time. \interp

\paragraph{The activation function.} The accumulated decoherence is measured by the activation function $E(a)=\exp(1-1/a)$ (Eq.~\eqref{eq:vc_E}), where $a$ is
the cosmological scale factor normalized to unity today. Its properties follow directly from the physical identification: $E\to0$ as $a\to0$ (negligible at
early times, $E(z{=}10)=4.5\times10^{-5}$); $E(1)=1$ today; $dE/da>0$ always (monotonically increasing, since each decoherence event is irreversible);
$E(\infty)=e$, the asymptotic limit, never reached. \calc

\paragraph{An arrow of time from irreversible records.} The arrow of time in this framework is not statistical. It does not depend on special initial
conditions or on the particular entropy state at the Big Bang. It is thermodynamically necessary: each quantum-to-classical transition writes one entry
permanently into the cosmic ledger at Landauer cost. The ledger only accumulates. $E(a)$ is monotonically increasing because each entry is irreversible.
This is the arrow of time obtained from the irreversibility of individual decoherence events, not imposed from boundary conditions. \interp

The sector separation follows immediately. Photons on null geodesics do not accumulate proper time and do not participate in gravitational decoherence.
They do not write to the ledger. The informational modification to gravitational dynamics applies exclusively to the matter sector: $\mu(a)<1$ with
$\Sigma(a)=1$ (Eq.~\eqref{part2:eq:mu_virial}), in the standard modified-growth parameters of Bertschinger \& Zukin~\cite{BertschingerZukin2008}. The photon
sector is geometrically standard. The matter sector carries the accumulated thermodynamic history of 13.8 billion years of decoherence transactions. \derived

\section{The black hole as encoding vault}\label{sec:vd_vault}
\paragraph{Why structure would require a local encoding surface.} The virial partition, on this reading, cannot close on one half alone. The geometric half
requires the informational half to close simultaneously, because the potential half stabilizes in spacetime curvature only when the kinetic half is absorbed
by an available encoding surface. If the informational half is produced but has no surface to write to, the thermodynamic accounting cannot balance, and the
structure cannot hold. \conjecture

The cosmological apparent horizon is always available, but it is cold ($T_H\approx2.65\times10^{-30}$\,K today, $H_0=67.16$) and expands uniformly, indifferent to
individual localities. For rapid local collapses on dynamical timescales $t_{\rm dyn}\ll H^{-1}$, the proposal is that the cosmic horizon cannot absorb the
locally produced information fast enough, so a local hot encoding surface is required. That premise is open: in IAM the horizon is always available, and the rate at which a local region must write is not set by $E(a)$ (Chapter~\ref{ch:satellites}). \openprob

The black hole horizon is the candidate surface. It forms when the local information production from gravitational decoherence saturates the holographic
bound on the bounding surface,
\begin{equation}
\frac{S_{\rm info}(R)}{A(R)}\ge\frac{k_B}{4\ell_P^2}\qquad(\text{one nat per }4\ell_P^2),\label{eq:vd_sat}
\end{equation}
(Chapter~\ref{ch:blackholes}). \conjecture\ Once formed, the black hole absorbs the informational half of the local virial partition permanently. Its
encoding throughput, the Hawking power of a Schwarzschild hole divided by $k_BT_{\rm BH}\ln2$ per bit at its own Hawking temperature $T_{\rm BH}=\hbar c^3/8\pi GMk_B$, is
\begin{equation}
\Gamma_{\rm BH}=\frac{P_H}{k_BT_{\rm BH}\ln2}=\frac{c^3}{1920\,GM\ln2}\ \text{bits\,s}^{-1},\qquad P_H=\frac{\hbar c^6}{15360\pi G^2M^2},\label{eq:vd_gamma}
\end{equation}
\derived\ (sympy, \texttt{verify\_virial\_papers.py} section E). The black hole is not a pipe to the cosmic horizon --- Hawking evaporation timescales
($2.1\times10^{67}$\,yr for a solar-mass black hole \calc) render it informationally negligible on cosmological timescales. It is a permanent local vault.


\section{Discussion}\label{sec:vd_discussion}
\paragraph{Time as the accumulated cost of classical existence.} The framework addresses a question that has resisted purely geometric treatment: why does
time have a direction? The answer proposed is not statistical. It is thermodynamic. Each quantum-to-classical transition pays an irreversible Landauer cost
to the nearest available encoding surface. The accumulated total of these costs, written on the cosmic horizon, is $E(a)$. Time has a direction because each
decoherence event is irreversible. The ledger only grows. \interp

This is consistent with the programme of temporal naturalism (Smolin 2013~\cite{Smolin2013}) in that it provides a physical mechanism for the
irreversibility of time. It extends that programme in one direction: the arrow of time is not a fundamental postulate but a consequence of the
thermodynamics of information production at every quantum-to-classical boundary. It should be noted that Smolin's temporal naturalism includes a stronger
thesis: not only that time is real, but that the laws of physics themselves evolve in time. IAM does not make this claim. The thermodynamic laws governing
$E(a)$ --- the first law, the second law, Landauer's principle --- are taken as fixed. The arrow of time emerges from the irreversibility of processes governed
by those laws, not from the evolution of the laws themselves. Whether the deeper connection between IAM's irreversibility and Smolin's evolving-laws
programme constitutes a tension or a complementarity is an open question not resolved here. \openprob\ $E(a)$ is not a parameter imposed from outside the
theory: it is fixed by horizon thermodynamics and the measured expansion history.

\paragraph{Dissipation-driven adaptation and the closure condition.} England's dissipation-driven adaptation (England 2013~\cite{England2013},
2015~\cite{England2015}) establishes a general thermodynamic principle: systems subject to a thermodynamic drive tend, over time, toward configurations that
more efficiently absorb and dissipate energy from their environment. Systems that align with the thermodynamic flow endure; those that cannot dissipate
sufficiently disperse. The probability of a driven system ending up in a given macroscopic configuration is proportional to the thermodynamic flux that
occurs on the way to that configuration.

The virial partition closure condition of Section~\ref{sec:vd_vault} is proposed as a specific gravitational instance of this principle. A collapsing dark
matter halo is a driven system: gravitational potential energy is being converted irreversibly into kinetic degrees of freedom at the virialization
boundary. On the proposal, the system endures as a galaxy if and only if it can route the informational half of its virial partition to a local encoding
surface --- a black hole horizon. Where no such surface can form, the dissipation cannot complete, and the structure disperses. The closure condition would
then be a necessary condition for endurance in England's sense, obtained for gravitational $1/r$ systems. The connection is meant as more than an analogy: the
Landauer cost at the virialization boundary is the physical quantity England's framework identifies as the thermodynamic flux required for stable
self-organization, and IAM computes this cost explicitly for gravitational systems and names the encoding surface that absorbs it. \conjecture\ Whether the full England programme ---
which addresses driven non-equilibrium systems far more generally --- shares deeper mathematical structure with the IAM virial partition is an open question.
\openprob

\paragraph{Thermal time and accumulated decoherence.} Rovelli's thermal time hypothesis (Rovelli 1993~\cite{Rovelli1993}; Connes \& Rovelli
1994~\cite{ConnesRovelli1994}) proposes that time emerges as a statistical effect from the thermodynamic state of a system: what we call the flow of time is
the modular flow of a Gibbs state, and the direction of time is rooted in thermodynamic irreversibility rather than in any fundamental asymmetry of the laws
of physics. The more recent work of Rovelli (2019~\cite{Rovelli2019}, 2022~\cite{Rovelli2022}) asks specifically where the low entropy of the early universe
was located and how thermodynamic traces constitute physical memory.

The activation function $E(a)$ is consistent with this programme and extends it in a specific direction. Rather than time emerging from incomplete knowledge
of a statistical ensemble, $E(a)$ in IAM records the accumulated irreversible Landauer costs of individual quantum-to-classical transitions written
permanently on the cosmic horizon. These are not statistical: each transaction is a specific physical event with a specific thermodynamic cost $k_BT_H\ln2$.
It accumulates not because we lack information about the microstate but because each decoherence event is physically irreversible at the level of the
second law. The record has a specific physical carrier ($E(a)$), a specific beginning (the electroweak transition, after which massive particles exist; it grows as bound structure forms), a
specific asymptotic limit ($e$), and specific cosmological observational consequences ($\mu(a)<1$, $\sigma_8=0.800$, testable by Euclid and DESI).
\interp\ Whether this constitutes a derivation of a thermal time variable in Rovelli's sense, or a complementary but distinct account of temporal
irreversibility, is an open question. \openprob

\paragraph{Cosmological natural selection and the black hole role.} Smolin's cosmological natural selection (Smolin 1992~\cite{Smolin1992},
1997~\cite{Smolin1997}) proposes that black holes are reproductive units of baby universes, with constants selected by a natural-selection-like process across
generations. IAM assigns black holes a different but complementary role: they are proposed as local encoding vaults, the thermodynamic structures that allow
galaxies to persist by absorbing the informational half of their virial partitions. Both frameworks place black holes at the centre of cosmic structure.
Whether these roles are independent, complementary, or aspects of a deeper unification is an open problem the IAM framework does not claim to resolve; it is
noted as a direction worth examining. \conjecture

\section{Open problems}\label{sec:vd_open}
\begin{enumerate}
\item The local-to-global encoding criterion: the condition under which decoherence events are assigned to a local black-hole vault rather than to the
global cosmic horizon. \openprob
\item The growth suppression from $\mu<1$ is derived from $\beta_m=\Omega_m/2$: a $13.6\,\%$ reduction of the coupling
today, which lowers $f\sigma_8$ by $4.25\,\%$ at $z=0$ with the same early amplitude. No measurement has yet confirmed or excluded it
(Chapter~\ref{ch:virial_tests}). \prediction
\end{enumerate}

Each item above is a concrete project. A cosmologist can run the activation function against the growth data now arriving, and a physicist working on decoherence can take any one of the open items into the laboratory.
